% ============================================================
% Question Bank: ch03-07 Converting Rational Numbers to Decimals
% Topic Title: Converting Rational Numbers to Decimals
% CCSS: 7.NS.A.2.D
% Questions: 30 (18 MC + 7 SA + 5 Graphical)
% ============================================================

% --- Multiple Choice Questions (18) ---

\begin{practiceQuestion}{3-7-q01}{mc}
\begin{questionText}
What is $\dfrac{3}{8}$ as a decimal?
\end{questionText}
\choiceA{$0.38$}
\choiceB{$0.375$}
\choiceC{$0.3\overline{8}$}
\choiceD{$0.83$}
\correctAnswer{B}
\explanation{$3 \div 8 = 0.375$. This is a terminating decimal.}
\end{practiceQuestion}

\begin{practiceQuestion}{3-7-q02}{mc}
\begin{questionText}
What is $\dfrac{1}{3}$ as a decimal?
\end{questionText}
\choiceA{$0.3$}
\choiceB{$0.33$}
\choiceC{$0.\overline{3}$}
\choiceD{$0.13$}
\correctAnswer{C}
\explanation{$1 \div 3 = 0.333\ldots$ The digit $3$ repeats forever: $0.\overline{3}$.}
\end{practiceQuestion}

\begin{practiceQuestion}{3-7-q03}{mc}
\begin{questionText}
Which fraction converts to a terminating decimal?
\end{questionText}
\choiceA{$\dfrac{1}{3}$}
\choiceB{$\dfrac{5}{6}$}
\choiceC{$\dfrac{7}{8}$}
\choiceD{$\dfrac{2}{9}$}
\correctAnswer{C}
\explanation{$\frac{7}{8} = 0.875$. The denominator $8 = 2^3$ has only $2$ as a prime factor, so the decimal terminates.}
\end{practiceQuestion}

\begin{practiceQuestion}{3-7-q04}{mc}
\begin{questionText}
What is $\dfrac{5}{11}$ as a decimal?
\end{questionText}
\choiceA{$0.\overline{45}$}
\choiceB{$0.45$}
\choiceC{$0.\overline{5}$}
\choiceD{$0.511$}
\correctAnswer{A}
\explanation{$5 \div 11 = 0.454545\ldots$ The digits $45$ repeat: $0.\overline{45}$.}
\end{practiceQuestion}

\begin{practiceQuestion}{3-7-q05}{mc}
\begin{questionText}
What is $\dfrac{7}{20}$ as a decimal?
\end{questionText}
\choiceA{$0.72$}
\choiceB{$0.35$}
\choiceC{$0.307$}
\choiceD{$3.5$}
\correctAnswer{B}
\explanation{$7 \div 20 = 0.35$. Terminates because $20 = 2^2 \times 5$.}
\end{practiceQuestion}

\begin{practiceQuestion}{3-7-q06}{mc}
\begin{questionText}
What is $-\dfrac{7}{12}$ as a decimal?
\end{questionText}
\choiceA{$-0.58$}
\choiceB{$-0.58\overline{3}$}
\choiceC{$-0.\overline{583}$}
\choiceD{$-0.712$}
\correctAnswer{B}
\explanation{$7 \div 12 = 0.5833\ldots = 0.58\overline{3}$. Add the negative sign.}
\end{practiceQuestion}

\begin{practiceQuestion}{3-7-q07}{mc}
\begin{questionText}
Which decimal is equivalent to $\dfrac{4}{9}$?
\end{questionText}
\choiceA{$0.4$}
\choiceB{$0.49$}
\choiceC{$0.\overline{4}$}
\choiceD{$0.44$}
\correctAnswer{C}
\explanation{$4 \div 9 = 0.444\ldots = 0.\overline{4}$.}
\end{practiceQuestion}

\begin{practiceQuestion}{3-7-q08}{mc}
\begin{questionText}
A fraction has a denominator whose only prime factors are $2$ and $5$. What kind of decimal will it produce?
\end{questionText}
\choiceA{Repeating}
\choiceB{Terminating}
\choiceC{It depends on the numerator}
\choiceD{Neither — it won't be a decimal}
\correctAnswer{B}
\explanation{If the denominator's only prime factors are $2$ and $5$, the decimal always terminates.}
\end{practiceQuestion}

\begin{practiceQuestion}{3-7-q09}{mc}
\begin{questionText}
What is $\dfrac{5}{6}$ as a decimal?
\end{questionText}
\choiceA{$0.56$}
\choiceB{$0.83$}
\choiceC{$0.8\overline{3}$}
\choiceD{$0.\overline{83}$}
\correctAnswer{C}
\explanation{$5 \div 6 = 0.8333\ldots = 0.8\overline{3}$. The $8$ does not repeat; only the $3$ repeats.}
\end{practiceQuestion}

\begin{practiceQuestion}{3-7-q10}{mc}
\begin{questionText}
What is $\dfrac{3}{5}$ as a decimal?
\end{questionText}
\choiceA{$0.35$}
\choiceB{$0.6$}
\choiceC{$0.\overline{6}$}
\choiceD{$0.53$}
\correctAnswer{B}
\explanation{$3 \div 5 = 0.6$. Terminates because $5$ is a factor of $10$.}
\end{practiceQuestion}

\begin{practiceQuestion}{3-7-q11}{mc}
\begin{questionText}
Which fraction produces a repeating decimal?
\end{questionText}
\choiceA{$\dfrac{1}{4}$}
\choiceB{$\dfrac{3}{10}$}
\choiceC{$\dfrac{1}{8}$}
\choiceD{$\dfrac{2}{3}$}
\correctAnswer{D}
\explanation{$\frac{2}{3} = 0.\overline{6}$. The denominator $3$ has a prime factor other than $2$ or $5$, so it repeats.}
\end{practiceQuestion}

\begin{practiceQuestion}{3-7-q12}{mc}
\begin{questionText}
What is $\dfrac{1}{6}$ as a decimal?
\end{questionText}
\choiceA{$0.16$}
\choiceB{$0.1\overline{6}$}
\choiceC{$0.\overline{16}$}
\choiceD{$0.\overline{6}$}
\correctAnswer{B}
\explanation{$1 \div 6 = 0.1666\ldots = 0.1\overline{6}$. The $1$ does not repeat; the $6$ repeats.}
\end{practiceQuestion}

\begin{practiceQuestion}{3-7-q13}{mc}
\begin{questionText}
What is $2\dfrac{1}{3}$ as a decimal?
\end{questionText}
\choiceA{$2.3$}
\choiceB{$2.\overline{3}$}
\choiceC{$2.13$}
\choiceD{$0.2\overline{3}$}
\correctAnswer{B}
\explanation{$2\frac{1}{3} = 2 + \frac{1}{3} = 2 + 0.\overline{3} = 2.\overline{3}$.}
\end{practiceQuestion}

\begin{practiceQuestion}{3-7-q14}{mc}
\begin{questionText}
What is $\dfrac{11}{25}$ as a decimal?
\end{questionText}
\choiceA{$0.44$}
\choiceB{$0.\overline{44}$}
\choiceC{$0.1125$}
\choiceD{$4.4$}
\correctAnswer{A}
\explanation{$11 \div 25 = 0.44$. Terminates because $25 = 5^2$.}
\end{practiceQuestion}

\begin{practiceQuestion}{3-7-q15}{mc}
\begin{questionText}
The decimal $0.\overline{27}$ means:
\end{questionText}
\choiceA{$0.27$}
\choiceB{$0.2777\ldots$}
\choiceC{$0.272727\ldots$}
\choiceD{$0.27027027\ldots$}
\correctAnswer{C}
\explanation{The bar over $27$ means the digits $2$ and $7$ repeat: $0.272727\ldots$}
\end{practiceQuestion}

\begin{practiceQuestion}{3-7-q16}{mc}
\begin{questionText}
Which decimal is equivalent to $-\dfrac{5}{8}$?
\end{questionText}
\choiceA{$-0.58$}
\choiceB{$-0.625$}
\choiceC{$-0.\overline{625}$}
\choiceD{$-0.585$}
\correctAnswer{B}
\explanation{$5 \div 8 = 0.625$. With the negative sign: $-0.625$.}
\end{practiceQuestion}

\begin{practiceQuestion}{3-7-q17}{mc}
\begin{questionText}
What is $\dfrac{7}{9}$ as a decimal?
\end{questionText}
\choiceA{$0.7$}
\choiceB{$0.79$}
\choiceC{$0.\overline{7}$}
\choiceD{$0.7\overline{9}$}
\correctAnswer{C}
\explanation{$7 \div 9 = 0.777\ldots = 0.\overline{7}$.}
\end{practiceQuestion}

\begin{practiceQuestion}{3-7-q18}{mc}
\begin{questionText}
Which statement is TRUE about rational numbers?
\end{questionText}
\choiceA{All fractions produce terminating decimals}
\choiceB{Repeating decimals are not rational numbers}
\choiceC{Every rational number can be written as a terminating or repeating decimal}
\choiceD{Decimals that go on forever are always irrational}
\correctAnswer{C}
\explanation{Every rational number (fraction) produces a decimal that either terminates or repeats. Both types are rational.}
\end{practiceQuestion}

% --- Short Answer Questions (7) ---

\begin{practiceQuestion}{3-7-q19}{sa}
\begin{questionText}
Convert $\dfrac{7}{8}$ to a decimal.
\end{questionText}
\correctAnswer{$0.875$}
\explanation{$7 \div 8 = 0.875$. Terminates.}
\end{practiceQuestion}

\begin{practiceQuestion}{3-7-q20}{sa}
\begin{questionText}
Convert $\dfrac{2}{11}$ to a decimal.
\end{questionText}
\correctAnswer{$0.\overline{18}$}
\explanation{$2 \div 11 = 0.181818\ldots = 0.\overline{18}$.}
\end{practiceQuestion}

\begin{practiceQuestion}{3-7-q21}{sa}
\begin{questionText}
Is $\dfrac{9}{16}$ a terminating or repeating decimal? Write the decimal.
\end{questionText}
\correctAnswer{Terminating: $0.5625$}
\explanation{$16 = 2^4$, so it terminates. $9 \div 16 = 0.5625$.}
\end{practiceQuestion}

\begin{practiceQuestion}{3-7-q22}{sa}
\begin{questionText}
Convert $-\dfrac{5}{9}$ to a decimal.
\end{questionText}
\correctAnswer{$-0.\overline{5}$}
\explanation{$5 \div 9 = 0.555\ldots = 0.\overline{5}$. Add the negative: $-0.\overline{5}$.}
\end{practiceQuestion}

\begin{practiceQuestion}{3-7-q23}{sa}
\begin{questionText}
Convert $\dfrac{3}{40}$ to a decimal.
\end{questionText}
\correctAnswer{$0.075$}
\explanation{$3 \div 40 = 0.075$. $40 = 2^3 \times 5$, so it terminates.}
\end{practiceQuestion}

\begin{practiceQuestion}{3-7-q24}{sa}
\begin{questionText}
Convert $1\dfrac{5}{6}$ to a decimal.
\end{questionText}
\correctAnswer{$1.8\overline{3}$}
\explanation{$\frac{5}{6} = 0.8\overline{3}$. So $1\frac{5}{6} = 1.8\overline{3}$.}
\end{practiceQuestion}

\begin{practiceQuestion}{3-7-q25}{sa}
\begin{questionText}
Write the repeating digits for $\dfrac{1}{7}$ as a decimal.
\end{questionText}
\correctAnswer{$0.\overline{142857}$}
\explanation{$1 \div 7 = 0.142857142857\ldots$ The six-digit block $142857$ repeats.}
\end{practiceQuestion}

% --- Graphical Questions (3 GMC + 2 GSA) ---

\begin{practiceQuestion}{3-7-q26}{gmc}
\begin{questionText}
The number line shows three points. Which point represents $\dfrac{3}{8}$?

\begin{center}
\begin{tikzpicture}[xscale=6]
  \draw[thick,->] (-0.1,0) -- (1.1,0);
  \foreach \x in {0,0.25,0.5,0.75,1} \draw (\x,0.08) -- (\x,-0.08) node[below] {\footnotesize $\x$};
  \fill[funRed] (0.375,0) circle (1.5pt) node[above=6pt] {A};
  \fill[funBlue] (0.5,0) circle (1.5pt) node[above=6pt] {B};
  \fill[funGreen] (0.625,0) circle (1.5pt) node[above=6pt] {C};
\end{tikzpicture}
\end{center}
\end{questionText}
\choiceA{Point A}
\choiceB{Point B}
\choiceC{Point C}
\choiceD{None of these}
\correctAnswer{A}
\explanation{$\frac{3}{8} = 0.375$, which is between $0.25$ and $0.5$. Point A is at $0.375$.}
\end{practiceQuestion}

\begin{practiceQuestion}{3-7-q27}{gmc}
\begin{questionText}
Below is a long division to convert $\dfrac{5}{6}$. What repeating decimal does this produce?

\begin{center}
\begin{tabular}{r@{\hskip 2pt}l}
  & $0.833\ldots$ \\
\cline{2-2}
$6\;)$ & $5.000\ldots$ \\
  & $4\,8$ \\
\cline{2-2}
  & \phantom{0}$20$ \\
  & \phantom{0}$18$ \\
\cline{2-2}
  & \phantom{00}$20$ \\
  & \phantom{00}$18$ \\
\cline{2-2}
  & \phantom{000}$2\ldots$ \\
\end{tabular}
\end{center}
\end{questionText}
\choiceA{$0.\overline{83}$}
\choiceB{$0.8\overline{3}$}
\choiceC{$0.83$}
\choiceD{$0.\overline{833}$}
\correctAnswer{B}
\explanation{The division shows $0.8333\ldots$ where the $8$ appears once and $3$ repeats: $0.8\overline{3}$.}
\end{practiceQuestion}

\begin{practiceQuestion}{3-7-q28}{gmc}
\begin{questionText}
The table sorts fractions into terminating and repeating decimals. Which fraction is in the wrong column?

\begin{center}
\begin{tabular}{|c|c|}
\hline
\textbf{Terminating} & \textbf{Repeating} \\
\hline
$\dfrac{1}{4}$ & $\dfrac{1}{3}$ \\[6pt]
\hline
$\dfrac{3}{5}$ & $\dfrac{2}{7}$ \\[6pt]
\hline
$\dfrac{1}{6}$ & $\dfrac{5}{9}$ \\[6pt]
\hline
\end{tabular}
\end{center}
\end{questionText}
\choiceA{$\dfrac{1}{4}$}
\choiceB{$\dfrac{3}{5}$}
\choiceC{$\dfrac{1}{6}$}
\choiceD{$\dfrac{2}{7}$}
\correctAnswer{C}
\explanation{$\frac{1}{6} = 0.1\overline{6}$ is a repeating decimal, not terminating. It belongs in the Repeating column.}
\end{practiceQuestion}

\begin{practiceQuestion}{3-7-q29}{gsa}
\begin{questionText}
Use the number line to locate $\dfrac{2}{3}$ and $\dfrac{3}{4}$ as decimals and determine which is larger.

\begin{center}
\begin{tikzpicture}[xscale=8]
  \draw[thick,->] (0.5,0) -- (1.05,0);
  \foreach \x/\l in {0.5/$0.5$,0.6/$0.6$,0.7/$0.7$,0.8/$0.8$,0.9/$0.9$,1.0/$1.0$} {
    \draw (\x,0.08) -- (\x,-0.08) node[below] {\footnotesize \l};
  }
  \fill[funRed] (0.667,0) circle (1pt) node[above=6pt] {\footnotesize $\frac{2}{3}$};
  \fill[funBlue] (0.75,0) circle (1pt) node[above=6pt] {\footnotesize $\frac{3}{4}$};
\end{tikzpicture}
\end{center}

Convert both fractions to decimals and state which is larger.
\end{questionText}
\correctAnswer{$\dfrac{2}{3} = 0.\overline{6} \approx 0.667$; $\dfrac{3}{4} = 0.75$; $\dfrac{3}{4}$ is larger}
\explanation{$\frac{2}{3} \approx 0.667$ and $\frac{3}{4} = 0.75$. Since $0.75 > 0.667$, $\frac{3}{4}$ is larger.}
\end{practiceQuestion}

\begin{practiceQuestion}{3-7-q30}{gsa}
\begin{questionText}
Complete the table by converting each fraction to a decimal. Then state whether it terminates or repeats.

\begin{center}
\begin{tabular}{|c|c|c|}
\hline
\textbf{Fraction} & \textbf{Decimal} & \textbf{Type} \\
\hline
$\dfrac{3}{4}$ & \rule{2cm}{0.4pt} & \rule{2cm}{0.4pt} \\[6pt]
\hline
$\dfrac{1}{9}$ & \rule{2cm}{0.4pt} & \rule{2cm}{0.4pt} \\[6pt]
\hline
$\dfrac{7}{20}$ & \rule{2cm}{0.4pt} & \rule{2cm}{0.4pt} \\[6pt]
\hline
\end{tabular}
\end{center}
\end{questionText}
\correctAnswer{$\dfrac{3}{4} = 0.75$ (terminating); $\dfrac{1}{9} = 0.\overline{1}$ (repeating); $\dfrac{7}{20} = 0.35$ (terminating)}
\explanation{$\frac{3}{4} = 0.75$ terminates ($4 = 2^2$). $\frac{1}{9} = 0.\overline{1}$ repeats ($9 = 3^2$). $\frac{7}{20} = 0.35$ terminates ($20 = 2^2 \times 5$).}
\end{practiceQuestion}
